\section*{Chapitre \ref{ch:hf:retail-wholesaling} --- Wholesaling de détail}

\begin{solution}{exo:hf:retail-wholesaling:1}
Le demi-spread vaut un cent~; 30~\% de celui-ci font 0.3 cent~: l'achat est exécuté à 50.017~\$. La fourchette effective vaut $2\times(50.017-50.010)=1.4$ cents, soit 70~\% des
deux cents affichés.
\end{solution}

\begin{solution}{exo:hf:retail-wholesaling:2}
$1.25(1-\pi)=0.1+0.1+0.2$~: $\pi=0.68$.
\end{solution}

\begin{solution}{exo:hf:retail-wholesaling:3}
E/Q $=1-\pi=0.70$~: les deux fourchettes varient proportionnellement au demi-spread.
\end{solution}

\begin{solution}{exo:hf:retail-wholesaling:4}
Sans effet de troupeau, les achats et les ventes s'alternent au hasard et l'inventaire dérive lentement~; avec effet de troupeau, des séries d'achats ou de ventes poussent l'inventaire au-delà de la limite, et
l'excédent doit être couvert~: 12.9~\% du volume contre 0.9~\%.
\end{solution}

\begin{solution}{exo:hf:retail-wholesaling:5}
Dans le modèle, les ordres institutionnels portent 1.5 cent d'information par action, soit plus que le demi-spread~: les exécuter au cours affiché fait perdre 0.57 cent par
action avant toute amélioration.
\end{solution}

\begin{solution}{exo:hf:retail-wholesaling:6}
Des différences dans les types d'ordres, leur taille et leur moment d'envoi d'un courtier à l'autre, les accords de tarification que les wholesalers passent avec chaque courtier (par exemple des engagements
d'amélioration), et des différences dans le routage propre à chaque courtier.
\end{solution}

\begin{solution}{exo:hf:retail-wholesaling:7}
Le gain net moyen est maximal pour une limite de 10\,000 actions (0.45 cent)~; la moyenne diminuée de l'écart type est maximale à 5\,000 (0.39 moins 0.21).
\end{solution}

\begin{solution}{exo:hf:retail-wholesaling:8}
La E/Q mesure l'amélioration par rapport au cours affiché à l'arrivée de l'ordre, et non le prix que le client aurait pu obtenir ailleurs~; un même ordre peut recevoir des
prix différents selon le courtier, et le paiement pour flux d'ordres n'entre pas dans la statistique.
\end{solution}

\begin{solution}{pb:hf:retail-wholesaling:1}
\begin{enumerate}
\item $h(1-\pi)-a-p-c_H+I$~: capture, mark-out, paiement pour flux d'ordres, couverture, P\&L d'inventaire.
\item Voir la \cref{def:hf:retail-wholesaling:segmentation}.
\item Détail~: 20\,000 ordres, médiane de 50 actions, 0.1 cent d'information, effet de troupeau~; institutionnels~: médiane de 2\,000 actions, 1.5 cent.
\item L'amélioration pour le trader, le paiement au courtier et la marge du wholesaler.
\item Détail~: 0.85, 0.475 et $-0.025$ cent par action~; institutionnels~: $-0.57$, $-0.94$, $-1.44$.
\item À 68~\% du demi-spread.
\item Voir la \cref{def:hf:retail-wholesaling:rate}~; 87~\%.
\item Fourchette affichée de 2.5 cents, fourchette effective de 1.75, E/Q de 0.70, amélioration de 0.375 cent par action.
\item Coûts de couverture face au risque d'inventaire~: 0.20 net avec un écart type de 0.06 à 1\,000 actions, 0.45 avec 0.38 à 10\,000, dominé par l'inventaire à 50\,000
(écart type de 1.32).
\item Il rend le flux unilatéral par vagues~: 12.9~\% du volume couvert contre 0.9~\% sans lui.
\item La SEC a retiré en juin 2025 les propositions sur la concurrence entre ordres et sur la meilleure exécution~; les modifications de la Rule 605 s'appliquent désormais à partir d'août 2026~; le
\omterm{def:hf:retail-wholesaling:rlp}{programme de liquidité de détail} (retail liquidity programme) de la NYSE a été approuvé en 2012 avec une amélioration inférieure au cent.
\item Voir la \cref{def:hf:retail-wholesaling:rlp}.
\item Avec une amélioration de 30~\% et 0.1 cent payé pour le flux, le wholesaler conserve 0.475 cent par action avant le P\&L d'inventaire (0.44 après), et le trader
reçoit 0.375 cent par action pour une fourchette affichée de 2.5 cents (E/Q de 0.70).
\item Les trois à la fois~: le modèle répartit 0.95 cent par action entre 0.375, 0.1 et 0.475.
\item Les wholesalers ont tarifé différemment les mêmes ordres selon le courtier~; les publications devraient présenter les prix par rapport à des exécutions comparables, courtier par courtier.
\item Beaucoup d'ordres de détail auraient été exposés à des enchères où d'autres auraient pu rivaliser.
\item Internalisation du détail sur actions.
\item À partir des mark-outs du flux par courtier, ainsi que de la taille et du moment des ordres.
\item Ses statistiques d'amélioration et d'exécution sur les ordres propres du courtier, comparées à celles des autres wholesalers.
\item Le droit de négocier contre un flux qui en sait peu, et l'obligation de le tarifer correctement.
\end{enumerate}
\end{solution}

\begin{solution}{iq:hf:retail-wholesaling:1}
Le wholesaler sait que le flux provient du détail et qu'il est presque non informé~; les cotations des bourses doivent se protéger contre tout le monde, y compris les traders informés, et
sont donc plus larges.
\iqlookfor{segmentation.}
\end{solution}

\begin{solution}{iq:hf:retail-wholesaling:2}
Les mark-outs de ses ordres à plusieurs horizons, par taille, type et moment de l'ordre, comparés à un flux de référence.
\iqlookfor{mark-outs par segment.}
\end{solution}

\begin{solution}{iq:hf:retail-wholesaling:3}
Des cotations issues d'un flux consolidé, les conditions par courtier dans une table de correspondance, l'inventaire comme état partagé~; une fonction de ces trois éléments calculée sur chaque ordre,
avec les limites vérifiées dans le même chemin.
\iqlookfor{conditions précalculées, inventaire partagé.}
\end{solution}

\begin{solution}{iq:hf:retail-wholesaling:4}
Des limites d'inventaire et de perte par valeur~; une couverture au-delà sur les bourses~; une amélioration plus faible ou nulle sur la valeur, et un routage vers les bourses si
l'inventaire ne peut pas être couvert.
\iqlookfor{limites et retrait progressif.}
\end{solution}

\begin{solution}{iq:hf:retail-wholesaling:5}
À partir du gain net par action du flux de ce courtier après amélioration, couverture et inventaire~: ne payer que ce qui laisse une marge, et comparer avec
l'amélioration que reçoivent les clients du courtier.
\iqlookfor{économie par courtier.}
\end{solution}

\begin{solution}{iq:hf:retail-wholesaling:6}
Pour un processus de signes AR(1), la variance de la somme croît comme $N(1+\rho)/(1-\rho)$ pour $N$ grand~: la persistance multiplie la variance du déséquilibre.
\iqlookfor{variance de long terme d'un AR(1).}
\end{solution}
